\chapter{Introduction and Modeled Fragment}\label{chap:intro}

This formalization models a fragment of the AWS IAM policy evaluation
logic in Lean~4 and proves that certain policy transforms can only
\emph{narrow} the set of requests a policy allows, never widen it.

\section{Three-Valued Condition Semantics}

Condition evaluation uses a three-valued logic (\texttt{Tri}):
\textbf{T}~(satisfied), \textbf{F}~(not satisfied),
\textbf{U}~(unknown---key absent from context or unsupported operator).
The connectives follow Kleene semantics.

The resolution rule: \textbf{UNKNOWN resolves toward visibility}.
An Allow statement applies when its condition is T or U (grant-preserving).
A Deny statement applies only when its condition is T (conservative).

\section{What ``allows'' Over-Approximates}

The \texttt{allows} function returns \texttt{true} if the policy
grants the request under the given condition context.  When the context
is incomplete (\texttt{noContext}), unknown conditions resolve toward
granting.  This is a deliberate over-approximation: the no-context
report never hides a grant that some complete context would allow.

\section{Layer Assumptions}

The SCP layer assumes full AWS access when no Allow statement is present
at a level (\texttt{assumed-full-access} note).  RCP support is limited
to services in \texttt{data/rcp-services.txt}.

\section{Out of Scope}

The following are not modeled:
\begin{itemize}
\item Identity-side policies in cross-account access (deferred to a future slice).
\item Session policies.
\item Resource-based policy evaluation beyond the resource document kind.
\end{itemize}

An outside collaborator should read these caveats before interpreting
any theorem as a guarantee about the full AWS policy evaluation engine.
The theorems guarantee properties of \emph{this model}, which faithfully
represents the modeled fragment and over-approximates on the boundaries
listed above.
